\documentclass{article}
\usepackage[utf8]{inputenc}
\usepackage[czech]{babel}
\usepackage{graphicx}
\usepackage{float}
\usepackage{amsmath}

\title{Státnicové otázky - Umělá inteligence}
\begin{document}

\maketitle

\section{Otázka 19: Metody prohledávání stavového prostoru}
\textbf{K dodělání:} 
\begin{itemize}
    \item Z původních poznámek se sem nehodí nic. Je nutné vypracovat celou otázku od nuly.
    \item Chybí: formulace problému, stavový prostor, strategie prohledávání.
    \item Chybí: neinformované a informované metody (včetně heuristik).
    \item Chybí: kompletnost a optimálnost algoritmů.
    \item Chybí: příklady problémů v AI řešených těmito metodami.
\end{itemize}



\section{Otázka 20: Biologicky inspirované výpočty}
\textbf{K dodělání:}
\begin{itemize}
    \item Z původních poznámek se sem nehodí nic. Je nutné vypracovat celou otázku od nuly.
    \item Chybí: rozdělení a principy evolučních a rojových algoritmů (GA, DE, PSO, ACO, GP).
    \item Chybí: lokální optimalizační metody (HC, HC-12, SA).
    \item Chybí: reprezentace a kódování řešení, surrogate optimalizace.
    \item Chybí: teorie "no free lunch" a prokletí dimenzionality.
    \item Chybí: příklady problémů v AI.
\end{itemize}


\section{Otázka 21: Neuronové sítě – vybraná základní paradigmata}
\textbf{K dodělání:}
\begin{itemize}
    \item Chybí: architektury ADALINE a RBF.
    \item Chybí: koncept kapacity sítě (capacity).
    \item Chybí: specifické příklady problémů v AI řešených těmito základními sítěmi.
\end{itemize}

\subsection{Formální neuron}
Vnitřní potenciál neuronu $z$ lze vypočítat dvěma základními způsoby:
\begin{itemize}
    \item \textbf{Lineární:} $z = w_0 + \sum^n_{i=1}x_iw_i$ (kde $w_0$ je bias, $x_i$ jsou vstupy a $w_i$ jsou váhy)
    \item \textbf{Radiální:} $z = w_0 + \sqrt{\sum^n_{i=1}(x_i-w_i)^2}$
\end{itemize}

\subsection{Architektury}
\subsubsection{Perceptron}
\begin{itemize}
    \item \textbf{Topologie:} Jednovrstvá přímovazební síť.
    \item \textbf{Aktivační funkce:} Nespojitá přenosová funkce (skoková, s prahem nebo bez).
    \item \textbf{Učení:} Iterativní, s učitelem. Váhy se adaptují na základě odchylky výstupu.
    \item \textbf{Omezení:} Funguje pouze jako lineární klasifikátor, řeší jen lineárně separabilní problémy (nelze realizovat funkci XOR).
\end{itemize}

\subsubsection{MLP (Vícevrstvý perceptron)}
\begin{itemize}
    \item \textbf{Topologie:} Vícevrstvá přímovazební síť (minimálně 3 vrstvy: vstupní, skrytá, výstupní). Úplné propojení neuronů mezi sousedními vrstvami.
    \item \textbf{Aktivační funkce:} Spojité a diferencovatelné funkce, např. sigmoida ($f(z) = \frac{1}{1+e^{-\sigma \cdot z}}$) nebo tangenta hyperbolická ($f(z) = \frac{1-e^{-\sigma \cdot z}}{1+e^{-\sigma \cdot z}}$).
    \item \textbf{Aplikace:} Modelování nelineárních funkcí a složitá klasifikace.
\end{itemize}

\subsubsection{SOM (Kohonenova samoorganizační mapa)}
\begin{itemize}
    \item \textbf{Topologie:} Dvouvrstvá síť s definovaným okolím ve výstupní vrstvě.
    \item \textbf{Učení:} Bez učitele, soutěžní strategie (kompetice). Aktivní je pouze jeden vítězný neuron.
    \item \textbf{Princip:} Ve výstupní vrstvě se počítá vzdálenost $d$ předloženého vzoru od vah neuronu: $d_j = \sum^n_{i=1} (x_i - w_{ij}(t))^2$. Vítězný neuron a jeho okolí adaptují své váhy, aby se přiblížily vstupu.
\end{itemize}

\subsection{Učení (Algoritmus Backpropagation)}
Učení s učitelem. Iterativní gradientní algoritmus, který minimalizuje kvadráty chybové funkce $E$.
\begin{itemize}
    \item \textbf{Výpočet chyby:} Chyba $E$ je rozdíl mezi skutečným výstupem $y$ a požadovaným výstupem $d$. Chyba pro jeden vzor: $^sE = \frac{1}{2}\sum_{j = 1}^{m}(^sy_j - ^sd_j)^2$
    \item \textbf{Adaptace vah:} Váhy se přepočítávají pro minimalizaci celkové chyby. Nová váha: $w_{ij}(t+1) = w_{ij}(t) + \Delta w_{ij}$, kde změna váhy je $\Delta w_{ij} = -\alpha \cdot \frac{\partial E_c}{\partial w_{ij}}$ ($\alpha$ je koeficient učení).
\end{itemize}

\subsection{Klasifikace, regrese a problémy učení}
\begin{itemize}
    \item \textbf{Klasifikace:} Zařazení vstupních dat do tříd (např. detekce spamu).
    \item \textbf{Regrese:} Odhad spojité výstupní hodnoty na základě vstupu.
    \item \textbf{Overfitting (Přeučení):} Stav, kdy se neuronová síť přespříliš adaptuje na trénovací množinu (naučí se i šum) a ztrácí schopnost generalizace na nová data.
\end{itemize}



\section{Otázka 22: Neuronové sítě – vybraná DNN paradigmata}
\textbf{K dodělání:}
\begin{itemize}
    \item Chybí: architektury RNN, Encoder-Decoder, Transformer (přehledově).
    \item Chybí: ukázka implementace (Keras).
    \item Chybí: příklady problémů v AI řešených pomocí DNN.
\end{itemize}

\subsection{Principy hlubokého učení (Deep Learning)}
Využívá vícevrstvé neuronové sítě pro extrakci rysů ze surových dat. V nižších vrstvách identifikují hrany a geometrické tvary, ve vyšších vrstvách detekují komplexnější struktury (např. písmena nebo obličeje).

\subsection{CNN (Konvoluční neuronová síť)}
Typické využití pro zpracování obrazu. Místo klasického násobení matic využívá konvoluci.
\begin{itemize}
    \item \textbf{Konvoluční vrstva:} Konvoluční jádro (kernel) o velikosti např. 3x3 až 11x11 prochází vstupní maticí obrazu.
    \item \textbf{Výpočet velikosti výstupu:} Po projití konvolucí se velikost počítá jako $(N-F)/S + 1$, kde $N$ je rozměr vstupu, $F$ je velikost kernelu a $S$ je velikost kroku (stride).
    \item \textbf{Max Pooling:} Slouží k redukci dimenzí vstupního signálu bez ztráty informační hodnoty. Z vybrané oblasti (např. 2x2) se vezme pouze maximální hodnota.
    \item \textbf{Plně propojená vrstva (Fully-Connected):} Výsledek po extrakci rysů je na konci zpracován klasickou plně propojenou sítí pro finální klasifikaci.
\end{itemize}

\subsection{Aktivační funkce pro DNN}
\begin{itemize}
    \item \textbf{ReLU (Rectified Linear Unit):} Pokud je hodnota menší než 0, nastaví ji na 0. Kladné hodnoty ponechává. $f(x) = \max(0, x)$.
    \item \textbf{Leaky ReLU:} Založena na ReLU, ale záporné hodnoty nenastaví na nulu, nýbrž jim ponechá mírný sklon (násobí je malou konstantou).
    \item \textbf{Softmax:} Používá se pro klasifikaci do více tříd ve výstupní vrstvě. Vrací vektor pravděpodobností se součtem 1. Počítá se jako $softmax(z)_i = \frac{e^{z_i}}{\sum_{j = 1}^{N}e^{z_j}}$.
\end{itemize}

\subsection{Trénink DNN}
\begin{itemize}
    \item \textbf{Batch learning:} Váhy se neupravují po každém vzoru, ale změny se akumulují pro celou dávku (batch) dat. Nové váhy se spočítají průměrem těchto změn, což stabilizuje učení a zrychluje proces.
\end{itemize}

\end{document}